\chapter{Dashboard Evaluation Results} \label{appendix:evaluation} 

\section{Participants} 
Three operational stakeholders participated in the evaluation of the dashboard suite. 

\begin{table}[h] 
    \centering 
    \begin{tabular}{|p{3cm}|p{8cm}|} 
        \hline 
        \textbf{Participant} & \textbf{Role} \\ 
        \hline\hline
        P1 & Clinical Oncology Director \\ 
        P2 & Radiotherapy Operations Manager \\ 
        P3 & Lead Treatment Radiographer \\ 
        \hline 
    \end{tabular} 
    \caption{Evaluation Participants} 
\end{table}

\section{System Usability Scale Results} 
\begin{table}[h] 
    \centering 
    \begin{tabular}{|l|c|} 
        \hline 
        \textbf{Participant} & \textbf{SUS Score} \\ 
        \hline\hline 
        P1 & 95.0 \\ 
        P2 & 90.0 \\ 
        P3 & 95.0 \\ 
        \hline 
        Mean & 93.3 \\ 
        \hline 
    \end{tabular} 
    \caption{System Usability Scale (SUS) Results} 
\end{table} 

The dashboard suite achieved a mean SUS score of 93.3, indicating a very high level of perceived usability.

\section{Operational Value Assessment} 

Participants rated the following statements using a five-point Likert scale (1 = Strongly Disagree, 5 = Strongly Agree). 

\begin{table}[h] 
    \centering 
    \begin{tabular}{|p{9cm}|c|} 
        \hline 
        \textbf{Statement} & \textbf{Mean Score} \\
        \hline\hline
        The dashboard improves visibility of post-MDT pathway performance & 5.0 \\
        The dashboard helps identify patients at risk of target breaches & 5.0 \\ 
        Displaying capacity and demand information supports operational decision-making & 4.7 \\ 
        The dashboard reduces reliance on manual reporting processes & 5.0 \\ 
        I would incorporate this dashboard into routine operational management activities & 4.7 \\ 
        \hline 
    \end{tabular} 
    \caption{Operational Value Assessment} 
\end{table}


\section{Design Principle Validation} 

\begin{table}[H]
    \centering 
    \begin{tabular}{|p{9cm}|c|} 
        \hline 
        \textbf{Statement} & \textbf{Mean Score} \\ 
        \hline\hline
        The dashboard helps identify risks before breaches occur & 5.0 \\ 
        Capacity information is presented in a useful and actionable way & 5.0 \\ 
        The pathway stages reflect meaningful operational events & 4.7 \\ 
        Information is refreshed frequently enough for operational use & 4.3 \\ 
        I trust the dashboard metrics & 4.0 \\ 
        The dashboard supports my day-to-day operational tasks & 4.3 \\ 
        \hline 
    \end{tabular} 
    \caption{Design Principle Evaluation Results} 
\end{table}
\clearpage
\section{Qualitative Feedback} 
\begin{table}[h] 
    \centering 
    \begin{tabular}{|p{4cm}|p{10cm}|} 
        \hline 
        \textbf{Theme} & \textbf{Feedback} \\ 
        \hline\hline
        Most Useful Features & Waterfall visualisation of pathway delays; filtering by oncologist and speciality; patient worklists and breach monitoring. \\ 
        Missing Information & SACT data; clinic utilisation; visibility of emergency inpatient referrals to radiotherapy; utilisation breakdown by machine availability versus service demand. \\ 
        Operational Decisions Supported & Referral monitoring; patient tracking; bottleneck identification; overtime planning; shift extensions; additional linear accelerator capacity; capacity planning based on expected referrals. \\
        Future Enhancements & Clinic-capacity modelling; CT-to-treatment decomposition; utilisation analysis; integration of structured electronic referrals and SACT datasets. \\ 
        General Comments & ``Very good. Can't wait to start using it.'' \newline ``Looks good.'' \newline Electronic referral implementation was identified as a useful future data source for before-and-after evaluation studies. \\ 
        \hline 
    \end{tabular} 
    \caption{Summary of Qualitative Feedback} 
\end{table}